\subsection{Volumeninterpretation}
\label{subsec:vector_algebra_scalar_triple_product_volume}

Die Vektoren $\vec a$ und $\vec b$ spannen ein Parallelogramm auf. Wird dieses
entlang von $\vec c$ verschoben, entsteht ein Spat. Wir möchten nun das Volumen
dieses Spats berechnen.

\noindent
\begin{minipage}[c]{0.57\linewidth}
  \[
    \begin{aligned}
      V
      &=A h \\
      &=|\vec a\times\vec b|\,|\vec c|\,|\cos(\varphi)|\\
      &=\bigl|(\vec a\times\vec b)\cdot\vec c\bigr|.
    \end{aligned}
  \]
\end{minipage}\hfill
\begin{minipage}[c]{0.39\linewidth}
  \centering
  \CachedTikzFigure{content/resources/chapters/vector_algebra/tikz/scalar_triple_product_volume.tex}
\end{minipage}

\begin{DefinitionBox}{Spatprodukt}{vector_algebra_scalar_triple_product}
  Das \emph{Spatprodukt} \EnglishTerm{scalar triple product} der Vektoren
  $\vec a,\vec b,\vec c\in\R^3$ ist der Skalar
  \[
    S=(\vec a\times\vec b)\cdot\vec c.
  \]
  Es ist ein orientiertes Volumen: $|S|$ ist das Volumen des aufgespannten
  Spats, das Vorzeichen von $S$ beschreibt seine Orientierung.
\end{DefinitionBox}

\begin{ExampleBox}{Volumen eines Spats}{vector_algebra_scalar_triple_product_volume_example}
  Die Vektoren
  \[
    \vec a=\begin{pmatrix}4\\2\\1\end{pmatrix},\qquad
    \vec b=\begin{pmatrix}3\\-1\\2\end{pmatrix},\qquad
    \vec c=\begin{pmatrix}0\\2\\-5\end{pmatrix}
  \]
  spannen einen Spat auf. Zunächst berechnen wir
  \[
    \vec a\times\vec b
    =\begin{ClozeContinuation}[19]{3}
      \begin{pmatrix}
        2\cdot2-1\cdot(-1)\\
        1\cdot3-4\cdot2\\
        4\cdot(-1)-2\cdot3
      \end{pmatrix}
      =\begin{pmatrix}5\\-5\\-10\end{pmatrix}
    \end{ClozeContinuation}.
  \]
  Damit beträgt das Volumen
  \[
    V
    =\begin{ClozeContinuation}[20]{3}
      \left|
        \begin{pmatrix}5\\-5\\-10\end{pmatrix}
        \cdot
        \begin{pmatrix}0\\2\\-5\end{pmatrix}
      \right|
      =|0-10+50|=40
    \end{ClozeContinuation}.
  \]
\end{ExampleBox}
